% -------------------------------------------------------------------------
% WBS ID: RES-MOD-SPEC
% Deliverable: Integrated Mathematical Model Specification
% Status: Draft complete
% Evidence status: Regional 2D NLSWE and local 3D URANS--VOF research specifications complete
% Review status: Not reviewed
% Last updated: 2026-07-18
% Dependencies: RES-MOD-EQS, RES-MOD-SRC, RES-MOD-IBC, RES-MOD-MUL, RES-PHY-SYN, RES-DAT-GEO, RES-DAT-SRC
% Downstream interfaces: RES-NUM-SPEC, RES-VER-SPEC, Software implementation
% -------------------------------------------------------------------------

\section{RES-MOD-SPEC --- Integrated Mathematical Model Specification}
\label{sec:res-mod-spec}

\subsection{Purpose and Scope}
\label{sec:res-mod-spec-purpose}

This Deliverable consolidates the continuous mathematical model for the
regional tsunami solver and the local three-dimensional hydrodynamic
tsunami--barrier model. The regional 2D NLSWE mathematical model and
the local 3D URANS--VOF hydrodynamic model are complete at
research-specification level. Remaining work belongs to
geospatial/domain selection, validation, implementation, detailed
coupling operators and activated FSI or structural-failure extensions.

\subsection{Research Questions}
\label{sec:res-mod-spec-research-questions}

The regional-model gate answers:

\begin{enumerate}
  \item Which horizontal governing model represents regional tsunami
    propagation and inundation?
  \item Which state variables, bed convention, source terms and
    condition roles define the mathematical model?
  \item Which decisions are established and which downstream
    workstreams remain open?
  \item Which continuous local model is used for pressure, force,
    moment, overtopping and barrier-performance comparison?
  \item Which local physics are baseline assumptions and which are
    deferred extensions?
\end{enumerate}

\subsection{Terminology and Definitions}
\label{sec:res-mod-spec-definitions}

% TODO: Define the technical terminology, symbols, quantities and
% classifications used in this Deliverable.

\subsection{Literature Evidence}
\label{sec:res-mod-spec-literature-evidence}

% TODO: Record evidence from peer-reviewed papers, authoritative datasets,
% standards, technical reports and primary documentation.
%
% Do not create a detached literature-summary list. Explain how each source
% supports, contradicts or limits a project decision.

\subsection{Governing Relations and Quantitative Definitions}
\label{sec:res-mod-spec-governing-relations}

% TODO: Record relevant equations, dimensionless groups, empirical models,
% extraction rules or quantitative definitions.
%
% Use align environments for displayed equations.
% Every displayed equation must have a unique \label{eq:...}.
% Do not place punctuation after displayed equations.

\subsection{System Boundary}
\label{sec:res-mod-spec-system-boundary}

The regional solver is a horizontal, depth-averaged,
cell-centred finite-volume model for the two-dimensional nonlinear
shallow-water equations. Its independent solved coordinates are
$(x,y,t)$. Vertical information is carried through scalar fields
$b(x,y,t)$, $\eta(x,y,t)$ and $h(x,y,t)=\eta-b$. Cross-sections are
diagnostic visualisations extracted from the horizontal solution rather
than separately solved cross-sectional models.

\subsection{State Variables and Parameters}
\label{sec:res-mod-spec-state-variables-parameters}

The conserved regional state is
\cref{eq:res-mod-eqs-regional-conserved-state}. The primitive state is
\cref{eq:res-mod-eqs-regional-primitive-state}. The bed convention is
positive upward, with $\eta=h+b$ and $h=\eta-b$.

\subsection{Continuous Governing Model}
\label{sec:res-mod-spec-continuous-governing-model}

The established regional governing model is the nonlinear
shallow-water balance law in
\cref{eq:res-mod-eqs-regional-nlswe-balance}, with Cartesian fluxes in
\cref{eq:res-mod-eqs-regional-x-flux,eq:res-mod-eqs-regional-y-flux}.
The model supports wetting and drying through the numerical method; it
does not introduce a regional volume-of-fluid phase variable.

The established local governing model is the three-dimensional
nonhydrostatic incompressible immiscible two-phase Navier--Stokes
system with VOF phase transport:
\cref{eq:res-mod-eqs-local-continuity,eq:res-mod-eqs-local-momentum,eq:res-mod-eqs-local-phase-transport}.
The local state is
\cref{eq:res-mod-eqs-local-fluid-state}; the operational turbulence
closure is URANS $k$--$\omega$ SST.

\subsection{Source Terms and Closures}
\label{sec:res-mod-spec-source-terms-closures}

The baseline regional source terms are bed slope, Manning friction and
open-boundary sponge damping. Dynamic earthquake forcing is represented
as a time-dependent bed field,
\cref{eq:res-mod-eqs-regional-dynamic-bed}. Since $h$ is the conserved
depth variable, no separate earthquake mass source is added to the
depth equation.

Coriolis and explicit lateral eddy viscosity remain disabled in the
baseline and may be activated only as later sensitivity studies.

The local effective viscosity is defined by
\cref{eq:res-mod-eqs-local-effective-viscosity}, with $\mu_t$ supplied
by the selected URANS $k$--$\omega$ SST closure in
\cref{eq:res-mod-src-sst-k-equation,eq:res-mod-src-sst-omega-equation,eq:res-mod-src-sst-eddy-viscosity}.
LES is retained only for high-fidelity comparison cases, not as the
baseline repeated-screening model.

\subsection{Initial, Boundary and Interface Conditions}
\label{sec:res-mod-spec-initial-boundary-interface-conditions}

The regional condition roles are those in
\cref{tab:res-mod-ibc-regional-boundary-roles}: characteristic or
radiation open boundaries for artificial ocean boundaries, optional
sponge damping for artificial cuts, coastline treatment by topography
and wetting--drying, reflective states for true walls, and conservative
regional-to-local transfer data for the later interface workstream.

The local condition roles are those in
\cref{tab:res-mod-ibc-local-boundary-roles}: regional-to-local inlet
forcing, open or radiation downstream and lateral boundaries, an open
atmospheric top, fixed no-slip terrain and a fixed no-slip barrier for
the hydrodynamic baseline.

\subsection{Model Coupling Requirements}
\label{sec:res-mod-spec-model-coupling-requirements}

The regional model supplies $h$, $\eta$, $q_x$, $q_y$, derived
velocities, boundary histories and checkpoint data to the local
3D/interface workstream. The local 3D model receives those histories
through the baseline one-way replay mode defined in `RES-MOD-MUL`.
Simultaneous coupling remains required when feedback or reflection
materially changes the quantities of interest.

\subsection{Sufficient Conditions for Validity}
\label{sec:res-mod-spec-sufficient-conditions-validity}

The regional model is valid only where the depth-averaged,
hydrostatic, long-wave assumptions are acceptable for the selected
outputs. It is not the model for local three-dimensional impact
pressure, vertical velocity structure, resolved overtopping jets, local
air--water interface topology or structural FSI.

The local model is valid as a baseline hydrodynamic model where the
barrier is fixed or weakly moving, impermeable or explicitly resolved,
and sufficiently ventilated that trapped-air compression is not
material. It is not a model for unresolved porous media, scour,
sediment transport, debris impact, structural failure or final design
certification without additional physics and validation.

\subsection{Candidate Approaches}
\label{sec:res-mod-spec-candidate-approaches}

% TODO: Define the credible candidate approaches considered.

\subsection{Critical Comparison}
\label{sec:res-mod-spec-critical-comparison}

% TODO: Compare candidates using physically and project-relevant criteria.
% Do not select an approach solely on popularity or implementation ease.

\subsection{Assumptions and Applicability}
\label{sec:res-mod-spec-assumptions}

% TODO: State assumptions and the conditions under which they are valid.

\subsection{Limitations and Failure Conditions}
\label{sec:res-mod-spec-limitations}

% TODO: State where the evidence, model or selected approach ceases to be
% reliable or applicable.

\subsection{Project-Specific Conclusions}
\label{sec:res-mod-spec-conclusions}

% TODO: Convert the evidence into conclusions specific to the Studentship
% project. Do not merely restate literature findings.

The regional 2D NLSWE formulation is now established at
research-specification level. It shall be treated as the initial
mathematical baseline for implementation and verification planning.
Future changes to the regional mathematical model are revisions to a
locked specification, not continuation of the initial model-selection
research.

The local 3D URANS--VOF formulation is also established at
research-specification level for hydrodynamic barrier-performance
assessment. Source roles: `3D-RANSVOF`, `3D-VOF`, `3D-SST`,
`3D-WAVEBC`, `3D-COUPLING` and `3D-DEFERRED`. Supporting sources are
assigned in the relevant Deliverables rather than repeated here as a
generic literature list.

\subsection{Selected Approach and Rationale}
\label{sec:res-mod-spec-selected-approach}

% TODO: Record the selected approach only after sufficient evidence exists.
% State the alternatives rejected and the reasons for rejection.

The selected regional mathematical approach is:

\begin{quote}
  horizontal two-dimensional, depth-averaged NLSWE with conserved state
  $[h,q_x,q_y]^{\mathsf{T}}$, positive-upward bathymetry, bed-slope and
  Manning sources, dynamic moving bed, open/radiation boundaries,
  sponge damping where required, and wetting--drying handled by the
  numerical method
\end{quote}

Dispersive extensions, regional-to-local interface algorithms and FSI
remain outside this regional mathematical-model lock.

The selected local mathematical approach is:

\begin{quote}
  three-dimensional incompressible immiscible two-phase
  Navier--Stokes equations, VOF water-volume fraction, mixture
  density and viscosity, URANS $k$--$\omega$ SST effective viscosity,
  regional-to-local inlet forcing, open/radiation artificial
  boundaries, fixed no-slip terrain and barrier surfaces, and
  hydrodynamic extraction of pressure, shear, force, moment, run-up,
  overtopping, transmission and energy metrics
\end{quote}

Porosity, trapped-air compressibility, scour, sediment transport,
deformable barriers, structural failure and final certification remain
deferred or external to the baseline hydrodynamic model.

\subsection{Unresolved Decisions}
\label{sec:res-mod-spec-unresolved-decisions}

% TODO: Record unresolved questions, required evidence and decision owners.

\begin{table}[htbp]
  \centering
  \small
  \renewcommand{\arraystretch}{1.2}
  \begin{tabular}{
      p{0.36\textwidth}
      p{0.46\textwidth}
    }
    \hline
    Established
    &
    Deferred or separate
    \\
    \hline
    Regional 2D NLSWE formulation; conserved and primitive states; bed
    convention; source-term catalogue; dynamic moving bed; open,
    reflective and sponge boundary roles
    &
    QGIS/geospatial domain lock; exact bathymetry/topography dataset;
    exact fault-source dataset; validation execution; C++
    implementation; local coupling operators; FSI model; optimisation
    and ML workflow
    \\
    \hline
  \end{tabular}
  \caption{Regional mathematical-model established/deferred status}
  \label{tab:res-mod-spec-regional-established-deferred}
\end{table}

\begin{table}[htbp]
  \centering
  \small
  \renewcommand{\arraystretch}{1.2}
  \begin{tabular}{
      p{0.36\textwidth}
      p{0.46\textwidth}
    }
    \hline
    Established
    &
    Deferred or separate
    \\
    \hline
    Local 3D incompressible two-phase Navier--Stokes model; VOF phase
    fraction; mixture properties; URANS $k$--$\omega$ SST baseline;
    local boundary partition; one-way replay role; pressure, shear,
    force, moment, run-up, overtopping and energy outputs
    &
    Exact VOF advection algorithm; pressure-correction sequence;
    numerical wall functions; mesh refinement; validation execution;
    simultaneous-coupling activation tolerance; FSI, structural
    failure, scour, sediment, detailed aeration and certification
    \\
    \hline
  \end{tabular}
  \caption{Local three-dimensional mathematical-model established/deferred status}
  \label{tab:res-mod-spec-local-established-deferred}
\end{table}

\subsection{Requirements and Handoffs}
\label{sec:res-mod-spec-requirements}

% TODO: State requirements passed to other Research Domains, Software,
% verification activities or later execution work.

`RES-NUM-SPEC` shall carry the corresponding discrete finite-volume
specification. `RES-DAT-GEO` and `RES-DAT-SRC` shall close the domain
and dataset selections. `RES-VER-SPEC` shall define the lake-at-rest,
wet/dry, moving-bottom, boundary and Tōhoku validation tests. Software
work shall not treat this document as evidence that implementation or
validation is complete.

For the local model, `RES-NUM-SPEC` shall carry the URANS--VOF
finite-volume solving sequence and convergence checks. `RES-VER-FRM`,
`RES-VER-BMK` and `RES-VER-MET` shall validate solver capability,
event-conditioned hydrodynamics and local impact/loading separately.

\subsection{Deliverable Completion Record}
\label{sec:res-mod-spec-completion-record}

% TODO: Record whether the Deliverable is:
%
% - Not started
% - In progress
% - Draft complete
% - Reviewed
% - Accepted
%
% Acceptance requires:
%
% - the research questions to be answered;
% - claims to be supported by appropriate evidence;
% - assumptions and limitations to be explicit;
% - the project-specific conclusion to be stated;
% - downstream requirements to be traceable;
% - unresolved decisions to be closed or formally deferred.

Regional and local hydrodynamic model status: draft complete for
research-specification purposes; review, validation and implementation
remain pending.
